% ============================================================
%  Chapitre 2 — Algèbre linéaire pour l'analyse des données
% ============================================================
\section{Algèbre Linéaire pour l'Analyse des Données}

% ────────────────────────────────────────────────────────────
\subsection{2.1 Rappels : valeurs et vecteurs propres}
% ────────────────────────────────────────────────────────────
\begin{frame}{2.1 — Rappels d'algèbre linéaire}
\begin{definition}[Valeur propre, vecteur propre]
Soit $\mathbf{A} \in \R^{p \times p}$. Un scalaire $\lambda \in \R$ est dit
\textbf{valeur propre} de $\mathbf{A}$ s'il existe $\mathbf{u} \in \R^p \setminus \{0\}$
tel que $\mathbf{A}\mathbf{u} = \lambda \mathbf{u}$. Le vecteur $\mathbf{u}$ est
le \textbf{vecteur propre} associé.
\end{definition}

\begin{columns}[T]
\column{0.49\textwidth}
\textbf{Polynôme caractéristique}
\[
P_{\mathbf{A}}(\lambda) = \det(\mathbf{A} - \lambda \mathbf{I}_p)
\]
Les valeurs propres sont les racines de $P_{\mathbf{A}}$ dans $\mathbb{C}$.
\vspace{0.2cm}

\textbf{Spectre} : $\sigma(\mathbf{A}) = \{\lambda_1, \ldots, \lambda_p\}$
(comptées avec multiplicité).

\textbf{Sous-espace propre}
\[
E_\lambda = \ker(\mathbf{A} - \lambda \mathbf{I}_p)
\]

\column{0.49\textwidth}
\begin{proposition}[Trace et déterminant]
\[
\Tr(\mathbf{A}) = \sum_{i=1}^p \lambda_i, \qquad
\det(\mathbf{A}) = \prod_{i=1}^p \lambda_i.
\]
\end{proposition}

\begin{proposition}[Diagonalisation]
$\mathbf{A}$ est \alert{diagonalisable} si et seulement si la somme des
dimensions des sous-espaces propres vaut $p$ :
\[
\sum_{\lambda \in \sigma(\mathbf{A})} \dim(E_\lambda) = p.
\]
On peut alors écrire $\mathbf{A} = \mathbf{P}\mathbf{D}\mathbf{P}^{-1}$
avec $\mathbf{D}$ diagonale.
\end{proposition}
\end{columns}
\end{frame}

% ────────────────────────────────────────────────────────────
\subsection{2.2 Théorème spectral}
% ────────────────────────────────────────────────────────────
\begin{frame}{2.2 — Théorème spectral pour matrices symétriques}
\begin{theorem}[Théorème spectral]
Soit $\mathbf{A} \in \R^{p \times p}$ une matrice \alert{symétrique réelle}
($\mathbf{A}^\top = \mathbf{A}$). Alors :
\begin{enumerate}
  \item Toutes les valeurs propres de $\mathbf{A}$ sont \textbf{réelles}.
  \item Les sous-espaces propres associés à des valeurs propres distinctes
        sont \textbf{orthogonaux}.
  \item Il existe une \textbf{base orthonormée} $\{\mathbf{u}_1, \ldots, \mathbf{u}_p\}$
        de $\R^p$ formée de vecteurs propres de $\mathbf{A}$.
\end{enumerate}
On a la décomposition :
\[
\boxed{\;\mathbf{A} = \mathbf{U}\, \mathbf{\Lambda}\, \mathbf{U}^\top
       = \sum_{k=1}^p \lambda_k \mathbf{u}_k \mathbf{u}_k^\top \;}
\]
avec $\mathbf{U}$ orthogonale ($\mathbf{U}^\top \mathbf{U} = \mathbf{I}_p$) et
$\mathbf{\Lambda} = \text{diag}(\lambda_1, \ldots, \lambda_p)$.
\end{theorem}

\begin{remark}[Conséquence majeure]
\small La matrice de variance-covariance $\mathbf{V}$ et la matrice de corrélation
$\mathbf{R}$ sont symétriques réelles : elles admettent toujours une base
\alert{orthonormée} de vecteurs propres. C'est le \textbf{cœur} de l'ACP.
\end{remark}
\end{frame}

% ────────────────────────────────────────────────────────────
\subsection{2.3 Formes quadratiques}
% ────────────────────────────────────────────────────────────
\begin{frame}{2.3 — Formes quadratiques et matrices SDP}
\begin{definition}[Forme quadratique]
À toute matrice symétrique $\mathbf{A} \in \R^{p\times p}$ on associe la
\textbf{forme quadratique}
\[
q_{\mathbf{A}} : \R^p \to \R, \quad
q_{\mathbf{A}}(\mathbf{x}) = \mathbf{x}^\top \mathbf{A}\, \mathbf{x}.
\]
\end{definition}

\begin{definition}[Matrice positive / définie positive]
\begin{itemize}
  \item $\mathbf{A}$ est \textbf{positive} ($\mathbf{A} \succeq 0$) si $q_{\mathbf{A}}(\mathbf{x}) \geq 0$ pour tout $\mathbf{x}$.
  \item $\mathbf{A}$ est \textbf{définie positive} ($\mathbf{A} \succ 0$) si $q_{\mathbf{A}}(\mathbf{x}) > 0$ pour tout $\mathbf{x} \neq 0$.
\end{itemize}
\end{definition}

\begin{theorem}[Caractérisation par les valeurs propres]
Soit $\mathbf{A}$ symétrique réelle. Alors :
\[
\mathbf{A} \succeq 0 \iff \forall k, \; \lambda_k \geq 0,
\qquad
\mathbf{A} \succ 0 \iff \forall k, \; \lambda_k > 0.
\]
\end{theorem}

\begin{remark}
$\mathbf{V}$ et $\mathbf{R}$ sont \alert{toujours} positives car
$\mathbf{V} = \frac{1}{n}\mathbf{X}_c^\top \mathbf{X}_c$ : pour tout $\mathbf{u}$,
$\mathbf{u}^\top \mathbf{V} \mathbf{u} = \frac{1}{n}\|\mathbf{X}_c \mathbf{u}\|^2 \geq 0$.
\end{remark}
\end{frame}

% ────────────────────────────────────────────────────────────
\subsection{2.4 Quotient de Rayleigh}
% ────────────────────────────────────────────────────────────
\begin{frame}{2.4 — Quotient de Rayleigh}
\begin{definition}[Quotient de Rayleigh]
Soit $\mathbf{A}$ symétrique réelle et $\mathbf{u} \in \R^p \setminus \{0\}$. Le
\textbf{quotient de Rayleigh} est :
\[
R_{\mathbf{A}}(\mathbf{u}) = \frac{\mathbf{u}^\top \mathbf{A}\, \mathbf{u}}{\mathbf{u}^\top \mathbf{u}}.
\]
\end{definition}

\begin{theorem}[Courant-Fischer-Weyl]
Soient $\lambda_1 \geq \lambda_2 \geq \cdots \geq \lambda_p$ les valeurs propres
de $\mathbf{A}$ (symétrique). Alors :
\[
\boxed{\;\lambda_1 = \max_{\mathbf{u} \neq 0} R_{\mathbf{A}}(\mathbf{u}),
\qquad
\lambda_p = \min_{\mathbf{u} \neq 0} R_{\mathbf{A}}(\mathbf{u}).\;}
\]
Plus généralement, pour le $k$-ième axe :
\[
\lambda_k = \max_{\substack{\mathbf{u} \perp \mathbf{u}_1, \ldots, \mathbf{u}_{k-1} \\ \mathbf{u} \neq 0}}
R_{\mathbf{A}}(\mathbf{u}).
\]
\end{theorem}

\begin{remark}[Importance pour l'ACP]
\small Maximiser la variance projetée $\mathbf{u}^\top \mathbf{V} \mathbf{u}$ sous
$\|\mathbf{u}\|=1$ revient à trouver le \alert{vecteur propre dominant} de $\mathbf{V}$.
C'est l'optimisation centrale de l'ACP.
\end{remark}
\end{frame}

% ────────────────────────────────────────────────────────────
\subsection{2.5 Décomposition en valeurs singulières}
% ────────────────────────────────────────────────────────────
\begin{frame}{2.5 — Décomposition en valeurs singulières (SVD)}
\begin{theorem}[SVD]
Toute matrice $\mathbf{X} \in \R^{n \times p}$ admet une décomposition
\[
\boxed{\;\mathbf{X} = \mathbf{U}\, \mathbf{\Sigma}\, \mathbf{V}^\top\;}
\]
où :
\begin{itemize}
  \item $\mathbf{U} \in \R^{n \times n}$ orthogonale ($\mathbf{U}^\top \mathbf{U} = \mathbf{I}_n$)
  \item $\mathbf{V} \in \R^{p \times p}$ orthogonale ($\mathbf{V}^\top \mathbf{V} = \mathbf{I}_p$)
  \item $\mathbf{\Sigma} \in \R^{n \times p}$ rectangulaire diagonale, de coefficients
        diagonaux $\sigma_1 \geq \sigma_2 \geq \cdots \geq \sigma_r \geq 0$,
        où $r = \rang(\mathbf{X})$.
\end{itemize}
Les $\sigma_k$ sont les \textbf{valeurs singulières}, les colonnes de
$\mathbf{U}$ et $\mathbf{V}$ sont les \textbf{vecteurs singuliers à gauche} et
\textbf{à droite}.
\end{theorem}

\begin{proposition}[Forme dyadique / SVD réduite]
\[
\mathbf{X} = \sum_{k=1}^r \sigma_k\, \mathbf{u}_k \mathbf{v}_k^\top.
\]
Chaque terme est une matrice de rang 1. La SVD décompose $\mathbf{X}$ en
\alert{somme} de matrices de rang 1, par ordre décroissant d'importance.
\end{proposition}
\end{frame}

% ────────────────────────────────────────────────────────────
\subsection{2.6 Lien spectral / SVD}
% ────────────────────────────────────────────────────────────
\begin{frame}{2.6 — Lien entre SVD et décomposition spectrale}
\begin{theorem}[Spectres de $\mathbf{X}^\top \mathbf{X}$ et $\mathbf{X} \mathbf{X}^\top$]
Soit $\mathbf{X} = \mathbf{U}\mathbf{\Sigma}\mathbf{V}^\top$ la SVD de $\mathbf{X}$. Alors :
\[
\mathbf{X}^\top \mathbf{X} = \mathbf{V}\, \mathbf{\Sigma}^\top \mathbf{\Sigma}\, \mathbf{V}^\top,
\qquad
\mathbf{X} \mathbf{X}^\top = \mathbf{U}\, \mathbf{\Sigma} \mathbf{\Sigma}^\top\, \mathbf{U}^\top.
\]
Les valeurs propres communes (non nulles) sont $\lambda_k = \sigma_k^2$.
\end{theorem}

\begin{columns}[T]
\column{0.49\textwidth}
\textbf{Démonstration (idée).}
\[
\mathbf{X}^\top \mathbf{X}
= (\mathbf{U}\mathbf{\Sigma}\mathbf{V}^\top)^\top (\mathbf{U}\mathbf{\Sigma}\mathbf{V}^\top)
= \mathbf{V}\mathbf{\Sigma}^\top \mathbf{U}^\top \mathbf{U} \mathbf{\Sigma}\mathbf{V}^\top
\]
$= \mathbf{V}\, (\mathbf{\Sigma}^\top \mathbf{\Sigma})\, \mathbf{V}^\top$
car $\mathbf{U}^\top \mathbf{U} = \mathbf{I}_n$. La matrice $\mathbf{\Sigma}^\top \mathbf{\Sigma}$
est diagonale de coefficients $\sigma_k^2$.

\column{0.49\textwidth}
\begin{block}{Conséquence pour l'ACP}
\small La SVD de $\mathbf{X}_c$ donne directement :
\begin{itemize}\small
  \item Les axes principaux : colonnes de $\mathbf{V}$
  \item Les composantes principales : $\mathbf{C} = \mathbf{X}_c \mathbf{V}$
  \item Les inerties expliquées : $\lambda_k = \sigma_k^2$
\end{itemize}
\end{block}

\begin{alertblock}{En pratique}
\small On utilise la \textbf{SVD numérique} pour calculer l'ACP, plus stable
que la diagonalisation directe de $\mathbf{V}$.
\end{alertblock}
\end{columns}
\end{frame}

% ────────────────────────────────────────────────────────────
\subsection{2.7 Théorème d'Eckart-Young}
% ────────────────────────────────────────────────────────────
\begin{frame}{2.7 — Théorème d'Eckart-Young-Mirsky}
\begin{theorem}[Eckart-Young, 1936]
Soit $\mathbf{X} = \sum_{k=1}^r \sigma_k \mathbf{u}_k \mathbf{v}_k^\top$ la SVD de $\mathbf{X}$.
Pour tout entier $q \leq r$, la matrice
\[
\mathbf{X}_q = \sum_{k=1}^q \sigma_k \, \mathbf{u}_k \mathbf{v}_k^\top
\]
est l'unique matrice de rang $\leq q$ minimisant l'erreur de Frobenius :
\[
\mathbf{X}_q = \argmin_{\rang(\mathbf{Y}) \leq q} \|\mathbf{X} - \mathbf{Y}\|_F^2,
\quad
\text{avec}
\quad
\|\mathbf{X} - \mathbf{X}_q\|_F^2 = \sum_{k=q+1}^{r} \sigma_k^2.
\]
\end{theorem}

\begin{block}{Norme de Frobenius}
\small $\displaystyle \|\mathbf{A}\|_F^2 = \sum_{i,j} a_{ij}^2 = \Tr(\mathbf{A}^\top \mathbf{A}) = \sum_k \sigma_k^2$.
\end{block}

\begin{remark}[Interprétation pour l'analyse des données]
\small Eckart-Young fonde \alert{toute la réduction de dimension} : la
\textbf{meilleure approximation de rang $q$} d'un nuage de points est la
projection sur le sous-espace engendré par les $q$ premiers vecteurs propres
de $\mathbf{X}_c^\top \mathbf{X}_c$ — c'est l'ACP.
\end{remark}
\end{frame}

% ────────────────────────────────────────────────────────────
\subsection{2.8 Projections orthogonales}
% ────────────────────────────────────────────────────────────
\begin{frame}{2.8 — Projections orthogonales}
\begin{definition}[Projecteur orthogonal]
Soit $F$ un sous-espace vectoriel de $\R^p$. Le \textbf{projecteur orthogonal}
sur $F$ est l'unique application linéaire $\mathbf{P}_F : \R^p \to F$ telle que
$\mathbf{x} - \mathbf{P}_F(\mathbf{x}) \perp F$ pour tout $\mathbf{x}$.
Il vérifie : $\mathbf{P}_F^2 = \mathbf{P}_F$ et $\mathbf{P}_F^\top = \mathbf{P}_F$.
\end{definition}

\begin{proposition}[Forme matricielle]
Si $\{\mathbf{u}_1, \ldots, \mathbf{u}_q\}$ est une base orthonormée de $F$ et
$\mathbf{U} = [\mathbf{u}_1 \,|\, \cdots \,|\, \mathbf{u}_q]$, alors :
\[
\mathbf{P}_F = \mathbf{U} \mathbf{U}^\top.
\]
La projection de $\mathbf{x}$ est $\mathbf{P}_F(\mathbf{x}) = \mathbf{U}(\mathbf{U}^\top \mathbf{x}) = \sum_{k=1}^q \langle \mathbf{x}, \mathbf{u}_k\rangle \mathbf{u}_k$.
\end{proposition}

\begin{theorem}[Théorème de Pythagore]
Pour tout $\mathbf{x} \in \R^p$ :
\[
\|\mathbf{x}\|^2 = \|\mathbf{P}_F(\mathbf{x})\|^2 + \|\mathbf{x} - \mathbf{P}_F(\mathbf{x})\|^2.
\]
\end{theorem}

\begin{remark}
\small En analyse des données : projeter sur l'espace propre des $q$ premiers axes
de $\mathbf{V}$ \alert{maximise} la norme projetée donc l'inertie projetée.
\end{remark}
\end{frame}

% ────────────────────────────────────────────────────────────
\subsection{2.9 Inertie projetée}
% ────────────────────────────────────────────────────────────
\begin{frame}{2.9 — Inertie projetée et meilleur sous-espace}
\begin{theorem}[Décomposition de l'inertie]
Soit $\mathcal{N}$ un nuage centré de matrice de variance-covariance $\mathbf{V}$,
et $F_q$ le sous-espace engendré par les $q$ premiers vecteurs propres de $\mathbf{V}$.
Alors l'inertie totale se décompose comme :
\[
\boxed{\;
I_{\text{totale}}
= \underbrace{\sum_{k=1}^q \lambda_k}_{\text{inertie projetée sur } F_q}
+ \underbrace{\sum_{k=q+1}^p \lambda_k}_{\text{inertie résiduelle}}.
\;}
\]
$F_q$ est le sous-espace de dimension $q$ \alert{maximisant} l'inertie projetée.
\end{theorem}

\begin{proposition}[Pourcentage d'inertie expliquée]
La part d'inertie expliquée par les $q$ premiers axes est :
\[
\tau_q = \frac{\lambda_1 + \cdots + \lambda_q}{\lambda_1 + \cdots + \lambda_p} \in [0, 1].
\]
On choisit $q$ pour que $\tau_q$ atteigne un seuil (ex : 80\%).
\end{proposition}
\end{frame}

% ────────────────────────────────────────────────────────────
\subsection{2.10 Métriques généralisées}
% ────────────────────────────────────────────────────────────
\begin{frame}{2.10 — Décomposition généralisée (avec métrique $\mathbf{M}$)}
\begin{theorem}[SVD généralisée]
Soient $\mathbf{M}$ une métrique sur $\R^p$ et $\mathbf{D}_p$ une métrique sur $\R^n$
(symétriques définies positives). Il existe une décomposition de $\mathbf{X}$
\[
\mathbf{X} = \mathbf{U}\, \mathbf{\Sigma}\, \mathbf{V}^\top
\]
avec $\mathbf{U}^\top \mathbf{D}_p \mathbf{U} = \mathbf{I}$ et $\mathbf{V}^\top \mathbf{M} \mathbf{V} = \mathbf{I}$.
Les vecteurs propres $\mathbf{v}_k$ diagonalisent $\mathbf{X}^\top \mathbf{D}_p \mathbf{X} \mathbf{M}$.
\end{theorem}

\begin{columns}[T]
\column{0.49\textwidth}
\textbf{Cadre unificateur}\\
\small Toutes les méthodes factorielles s'écrivent comme une SVD
généralisée d'un triplet $(\mathbf{X}, \mathbf{M}, \mathbf{D}_p)$ :

\begin{itemize}\small
  \item ACP : $\mathbf{M} = \mathbf{I}_p$ ou $\mathbf{D}_{1/s^2}$, $\mathbf{D}_p = \frac{1}{n}\mathbf{I}_n$
  \item AFC : $\mathbf{M} = \mathbf{D}_{1/x_{\bullet j}}$, $\mathbf{D}_p$ liée aux profils
  \item ACM : AFC du tableau disjonctif
  \item AFD : $\mathbf{M} = \mathbf{V}_{\text{intra}}^{-1}$
\end{itemize}
\column{0.49\textwidth}
\begin{block}{Vue duale}
\small Le \textbf{schéma de dualité} (Cailliez-Pagès) :
\[
\mathbf{X}^\top \mathbf{D}_p \mathbf{X} \mathbf{M}
\quad \stackrel{\text{spectres}}{\equiv} \quad
\mathbf{X} \mathbf{M} \mathbf{X}^\top \mathbf{D}_p
\]
Mêmes valeurs propres non nulles : c'est la dualité individus/variables.
\end{block}

\begin{alertblock}{Conséquence}
\small Une seule machine algébrique --- la SVD généralisée --- couvre
l'\emph{ensemble} des méthodes factorielles classiques.
\end{alertblock}
\end{columns}
\end{frame}

% ────────────────────────────────────────────────────────────
\subsection{2.11 Aperçu R}
% ────────────────────────────────────────────────────────────
\begin{frame}[fragile]{2.11 — Algèbre linéaire en R}
\begin{lstlisting}
# Matrice symetrique definie positive
set.seed(42)
A <- matrix(rnorm(16), 4, 4); A <- t(A) %*% A   # symetrique, SDP

# Decomposition spectrale
res <- eigen(A)
res$values    # valeurs propres (decroissantes)
res$vectors   # vecteurs propres en colonnes (orthonormes)

# Verification : A = U D U^T
U <- res$vectors; D <- diag(res$values)
all.equal(A, U %*% D %*% t(U))   # TRUE

# SVD d'une matrice generale
X <- matrix(rnorm(20), 5, 4)
svd_res <- svd(X)
svd_res$d     # valeurs singulieres
svd_res$u; svd_res$v

# Verification : X = U Sigma V^T
all.equal(X, svd_res$u %*% diag(svd_res$d) %*% t(svd_res$v))  # TRUE

# Approximation de rang 2 (Eckart-Young)
X2 <- svd_res$u[, 1:2] %*% diag(svd_res$d[1:2]) %*% t(svd_res$v[, 1:2])
norm(X - X2, type = "F")^2          # = somme des sigma^2 restants
sum(svd_res$d[3:4]^2)
\end{lstlisting}
\end{frame}

% ────────────────────────────────────────────────────────────
\subsection{2.12 Récapitulatif}
% ────────────────────────────────────────────────────────────
\begin{frame}{2.12 — Récapitulatif du chapitre}
\begin{columns}[T]
\column{0.49\textwidth}
\textbf{Théorèmes à retenir}
\begin{itemize}\small
  \item \textbf{Spectral} : $\mathbf{A}$ symétrique réelle
        $\Rightarrow$ base orthonormée de vecteurs propres
        $\mathbf{A} = \mathbf{U}\mathbf{\Lambda}\mathbf{U}^\top$
  \item \textbf{Courant-Fischer-Weyl} : $\lambda_1 = \max R_{\mathbf{A}}(\mathbf{u})$
  \item \textbf{SVD} : $\mathbf{X} = \mathbf{U}\mathbf{\Sigma}\mathbf{V}^\top$ pour toute matrice
  \item \textbf{Eckart-Young} : meilleure approximation de rang $q$
  \item \textbf{Pythagore} : $\|\mathbf{x}\|^2 = \|\mathbf{P}_F\mathbf{x}\|^2 + \|\mathbf{x} - \mathbf{P}_F\mathbf{x}\|^2$
\end{itemize}

\column{0.49\textwidth}
\begin{block}{Pont vers l'ACP}
Toutes les briques sont en place :
\begin{enumerate}\small
  \item Variance-covariance $\mathbf{V}$ symétrique SDP
  \item Théorème spectral $\Rightarrow$ axes propres orthogonaux
  \item Quotient de Rayleigh $\Rightarrow$ inertie projetée maximale
  \item Eckart-Young $\Rightarrow$ ACP est optimale
\end{enumerate}
\end{block}

\begin{alertblock}{Idée maîtresse}
L'ACP, l'AFC, l'ACM, l'AFD sont \alert{une seule et même opération} :
diagonaliser une matrice symétrique pour décomposer une inertie.
\end{alertblock}
\end{columns}
\end{frame}
